\SourcePage{371}{374}
\section{Relation between molecular and atomic terms}

As the distance between the nuclei of a diatomic molecule is increased, the limiting system consists of two isolated atoms (or ions). This raises the question of how an electronic term of the molecule corresponds to the atomic states obtained when the atoms are separated (E.~Wigner, E.~Witmer, 1928). The correspondence is not unique: bringing together two atoms in specified states may produce a molecule in several different electronic states.

First suppose that the molecule consists of two different atoms. Let the isolated atoms have orbital angular momenta $L_1$, $L_2$ and spins $S_1$, $S_2$, with $L_1\geqslant L_2$. The projections of the angular momenta on the line joining the nuclei take the values $M_1=-L_1,-L_1+1,\ldots,L_1$ and $M_2=-L_2,-L_2+1,\ldots,L_2$. When the atoms approach, the resulting $\Lambda$ is the magnitude of $M_1+M_2$. Combining every possible pair of $M_1$ and $M_2$, we find the following numbers of occurrences of the values $\Lambda=|M_1+M_2|$:
\[
\begin{array}{c@{\qquad}l}
\Lambda=L_1+L_2&2\text{ times},\\
L_1+L_2-1&4\text{ times},\\
\dotfill&\dotfill\\
L_1-L_2&2(2L_2+1)\text{ times},\\
L_1-L_2-1&2(2L_2+1)\text{ times},\\
\dotfill&\dotfill\\
1&2(2L_2+1)\text{ times},\\
0&2L_2+1\text{ times}.
\end{array}
\]
Remembering that every term with $\Lambda\ne0$ is doubly degenerate, whereas a term with $\Lambda=0$ is nondegenerate, we find that the possible terms are
\begin{equation}
\begin{array}{r@{\quad}c@{\quad}l}
1\text{ term}&\text{with}&\Lambda=L_1+L_2,\\
2\text{ terms}&\text{with}&\Lambda=L_1+L_2-1,\\
\dotfill&&\dotfill\\
2L_2+1\text{ terms}&\text{with}&\Lambda=L_1-L_2,\\
2L_2+1\text{ terms}&\text{with}&\Lambda=L_1-L_2-1,\\
\dotfill&&\dotfill\\
2L_2+1\text{ terms}&\text{with}&\Lambda=0;
\end{array}
\tag{80.1}
\end{equation}
altogether there are $(2L_2+1)(L_1+1)$ terms with $\Lambda$ ranging from 0 to $L_1+L_2$.

The two atomic spins $S_1$, $S_2$ combine into the total molecular spin according to the usual angular-momentum addition rule, giving the following

\SourcePage{372}{375}
possible values of $S$:
\begin{equation}
 S=S_1+S_2,\quad S_1+S_2-1,\ldots,|S_1-S_2|.
 \tag{80.2}
\end{equation}
Combining each of these values with all the values of $\Lambda$ in (80.1) gives the complete list of possible terms of the resulting molecule.

For $\Sigma$ terms their sign must also be determined. This is easily done by observing that, as $r\to\infty$, the molecular wave functions can be written as products, or sums of products, of the wave functions of the two atoms. The value $\Lambda=0$ can arise either from nonzero projections with $M_1=-M_2$ or from $M_1=M_2=0$. Denote the wave functions of the first and second atoms by $\psi_{M_1}^{(1)}$, $\psi_{M_2}^{(2)}$. For $M=|M_1|=|M_2|\ne0$, form the symmetric and antisymmetric products
\[
 \psi^+=\psi_M^{(1)}\psi_{-M}^{(2)}+\psi_{-M}^{(1)}\psi_M^{(2)},\qquad
 \psi^-=\psi_M^{(1)}\psi_{-M}^{(2)}-\psi_{-M}^{(1)}\psi_M^{(2)}.
\]
Reflection in a vertical plane containing the molecular axis reverses the angular-momentum projection on the axis, so that $\psi_M^{(1)}$, $\psi_M^{(2)}$ become $\psi_{-M}^{(1)}$, $\psi_{-M}^{(2)}$, respectively, and conversely. The function $\psi^+$ is then unchanged, while $\psi^-$ changes sign. The former therefore belongs to a $\Sigma^+$ term and the latter to a $\Sigma^-$ term. Thus each value of $M$ yields one term of each kind. Since $M$ has $L_2$ possible values ($M=1,\ldots,L_2$), there are in all $L_2$ terms of each sign.

If $M_1=M_2=0$, the molecular wave function has the form $\psi=\psi_0^{(1)}\psi_0^{(2)}$. To determine the behavior of $\psi_0^{(1)}$ under reflection in a vertical plane, choose coordinates centered on the first atom with the $z$ axis along the molecular axis. Reflection in the vertical $xz$ plane is equivalent to an inversion about the origin followed by a rotation through $180^\circ$ about the $y$ axis. Under inversion, $\psi_0^{(1)}$ acquires the factor $P_1$, where $P_1=\pm1$ is the parity of the state of the first atom. The effect of an infinitesimal, and hence of any finite, rotation on the wave function is completely fixed by the atom's total orbital angular momentum. It is therefore enough to consider the special case of a one-electron atom with orbital angular momentum $l$ and $z$ component $m=0$; replacing $l$ by $L$ in the result gives the answer for an arbitrary atom. Apart from a constant factor, the angular part of the electron wave function with $m=0$ is $P_l(\cos\theta)$ (see (28.8)). Rotation through $180^\circ$ about the $y$ axis gives $x\to-x$, $y\to y$, $z\to-z$, or, in spherical coordinates, $r\to r$, $\theta\to\pi-\theta$, $\varphi\to\pi-\varphi$. Consequently $\cos\theta\to-\cos\theta$, and $P_l(\cos\theta)$ acquires the factor $(-1)^l$.

\SourcePage{373}{376}
Thus reflection in a vertical plane multiplies $\psi_0^{(1)}$ by $(-1)^{L_1}P_1$. Similarly, $\psi_0^{(2)}$ is multiplied by $(-1)^{L_2}P_2$, and the product $\psi=\psi_0^{(1)}\psi_0^{(2)}$ acquires the factor $(-1)^{L_1+L_2}P_1P_2$. The term is $\Sigma^+$ or $\Sigma^-$ according as this factor is $+1$ or $-1$.

Collecting these results, we find that among the $(2L_2+1)$ $\Sigma$ terms of each possible multiplicity, $L_2+1$ are $\Sigma^+$ terms and $L_2$ are $\Sigma^-$ terms if $(-1)^{L_1+L_2}P_1P_2=+1$; the numbers are reversed if $(-1)^{L_1+L_2}P_1P_2=-1$.

We now turn to a molecule composed of identical atoms. The rules for adding the atomic spins and orbital angular momenta to form the molecular $S$ and $\Lambda$ remain the same as for unlike atoms. The new feature is that the terms may be even or odd. Two cases must be distinguished: the atoms may be in either different or identical states.

If the atoms are in different states\footnote{In particular, this may be the combination of a neutral atom with an ionized one.}, the total number of possible terms is doubled relative to the corresponding molecule of unlike atoms. Reflection through the origin, located at the midpoint of the molecular axis, interchanges the states of the two atoms. By symmetrizing or antisymmetrizing the molecular wave function under this interchange, we obtain two terms with the same $\Lambda$ and $S$, one even and the other odd. Hence equal numbers of even and odd terms are obtained.

If both atoms are in the same state, the total number of states is the same as for a molecule of unlike atoms. Their parities are given by the following results of an analysis that we omit because of its length\footnote{See E.~Wigner, E.~Witmer~// Zs. Physik. 1928. V.~51 S.~850.}.

\SourcePage{374}{377}
Let $N_g$ and $N_u$ be the numbers of even and odd terms having specified $\Lambda$ and $S$. Then
\[
\begin{array}{ll}
\text{if $\Lambda$ is odd,}&N_g=N_u,\\
\text{if $\Lambda$ is even and $S$ is even $(S=0,2,4,\ldots)$,}&N_g=N_u+1,\\
\text{if $\Lambda$ is even and $S$ is odd $(S=1,3,\ldots)$,}&N_u=N_g+1.
\end{array}
\]
Among the $\Sigma$ terms, $\Sigma^+$ and $\Sigma^-$ must finally be distinguished. The rule is
\[
\begin{array}{ll}
\text{if $S$ is even,}&N_g^+=N_u^-+1=L+1,\\
\text{if $S$ is odd,}&N_u^+=N_g^-+1=L+1
\end{array}
\]
(where $L_1=L_2\equiv L$). Every $\Sigma^+$ term has parity $(-1)^S$, and every $\Sigma^-$ term parity $(-1)^{S+1}$.

Besides the relation considered above between molecular terms and the atomic terms obtained as $r\to\infty$, one may also ask how the molecular terms are related to those of the ``united atom'' that would result as $r\to0$, when the two nuclei are brought to the same point; an example is the relation between terms of the $\mathrm H_2$ molecule and the He atom. The following rules are readily obtained. From a united-atom term with spin $S$, orbital angular momentum $L$, and parity $P$, separating the constituent atoms gives molecular terms with spin $S$ and axial angular momentum $\Lambda=0,1,\ldots,L$, one term for each value of $\Lambda$. The parity of the molecular term is the same as $P$ for the atomic term ($g$ for $P=+1$ and $u$ for $P=-1$). The molecular term with $\Lambda=0$ is a $\Sigma^+$ term if $(-1)^LP=+1$, and a $\Sigma^-$ term if $(-1)^LP=-1$.

\ProblemHeading 1. Determine the possible terms of the $\mathrm H_2$, $\mathrm N_2$, $\mathrm O_2$, and $\mathrm{Cl}_2$ molecules that can arise when atoms in their normal states combine.

\SolutionHeading The rules given above yield the following possible terms:
\[
\begin{array}{c@{\quad}l@{\quad}l}
\text{molecule }\mathrm H_2& (\text{atoms in }{}^2S\text{ states}):&{}^1\Sigma_g^+,\ {}^3\Sigma_u^+;\\
\mathrm N_2&(\text{atoms in }{}^4S\text{ states}):&{}^1\Sigma_g^+,\ {}^3\Sigma_u^+,\ {}^5\Sigma_g^+,\ {}^7\Sigma_u^+;\\
\mathrm{Cl}_2&(\text{atoms in }{}^2P\text{ states}):&2{}^1\Sigma_g^+,\ {}^1\Sigma_u^-,\ {}^1\Pi_g,\ {}^1\Pi_u,\ {}^1\Delta_g,\\
&&2{}^3\Sigma_u^+,\ {}^3\Sigma_g^-,\ {}^3\Pi_g,\ {}^3\Pi_u,\ {}^3\Delta_u;\\
\mathrm O_2&(\text{atoms in }{}^3P\text{ states}):&2{}^1\Sigma_g^+,\ {}^1\Sigma_u^-,\ {}^1\Pi_g,\ {}^1\Pi_u,\ {}^1\Delta_g,\\
&&2{}^3\Sigma_u^+,\ {}^3\Sigma_g^-,\ {}^3\Pi_u,\ {}^3\Pi_g,\ {}^3\Delta_u,\\
&&2{}^5\Sigma_g^+,\ {}^5\Sigma_u^-,\ {}^5\Pi_g,\ {}^5\Pi_u,\ {}^5\Delta_g.
\end{array}
\]
(A numeral placed before a term symbol gives the number of terms of that type when it exceeds one.)

\ProblemHeading 2. Do the same for HCl and CO.

\SourcePage{375}{378}
\SolutionHeading When unlike atoms combine, the parity of their states is also relevant. Formula (31.6) shows that the normal states of H, O, and C atoms are even, while that of Cl is odd (for the electron configurations of the atoms, see Table~3). The preceding rules give
\[
\begin{array}{c@{\quad}l@{\quad}l}
\text{molecule HCl}&(\text{atoms in }{}^2S_g,{}^2P_u\text{ states}):&{}^{1,3}\Sigma^+,\ {}^{1,3}\Pi;\\
\text{molecule CO}&(\text{atoms in }{}^3P_g,{}^3P_g\text{ states}):&2{}^{1,3,5}\Sigma^+,\ {}^{1,3,5}\Sigma^-,\\
&&2{}^{1,3,5}\Pi,\ {}^{1,3,5}\Delta.
\end{array}
\]
